\chapter{Dual-Core DIP Example}
\label{chap:fpgasteel-dualcore-dip}

\section{About this project}
\label{sec:fpgasteel-dualcore-dip-about}

Three sibling projects share this section, since they only differ in how the same algorithm is compiled, never in what it does: \cmd{dualcore_dip_example}, \cmd{dualcore_dip_example_no_neon}, and \cmd{dualcore_dip_example_no_neon_no_icache} (all under \repopath{sw/}). Paragraph~\ref{subsec:fpgasteel-dualcore-dip-build-variants} covers exactly what distinguishes the three.

All three take the same camera-to-\term{VGA} pipeline from \nameref{chap:fpgasteel-dualcore-forwarding} and replace the plain format conversion with a small computer-vision pipeline, \ccode{DualCoreDIP}: on every captured frame, it looks for one white ball against a black background and, if it finds one that is not touching the image border, draws its outline and a centroid cross on the \term{VGA} output. At a macroscopic level, the pipeline:

\begin{enumerate}
 \item converts the frame to grayscale (\term{RGB565} \cite{OV7670-DS} $\rightarrow$ \term{RGB24} + grayscale in the same pass, so both the display copy and the analysis copy are produced together);
 \item thresholds it to a binary image (a pixel counts as ``object'' above a fixed brightness);
 \item cleans the binary image up morphologically, an opening (erosion followed by dilation) to remove isolated noise pixels without changing the blob's overall size;
 \item rejects the frame early if the cleaned blob touches the image border (it cannot be fully measured) or if too few pixels survived thresholding (nothing there);
 \item extracts the blob's outline, bounding box, and centroid, and overlays all three on the \term{RGB24} frame the camera pipeline already produced.
\end{enumerate}

\repopath{main0.cpp} reports the outcome on the \term{LED}s (\term{LED0} = object found, \term{LED1} = touching the border, \term{LED2} = empty) and the per-frame processing time on the hex display, exactly as \nameref{chap:fpgasteel-dualcore-forwarding} reports its own conversion time.

Every stage above runs split across both cores, using the exact same \ccode{Worker}/\ccode{Core1} dispatch mechanism documented in Paragraph~\ref{subsec:fpgasteel-dualcore-forwarding-dispatch}: nothing new had to be invented here, this project is simply that mechanism invoked several times instead of once. Section~\ref{sec:fpgasteel-dualcore-dip-operating-principles} covers the pipeline and its cache discipline; the rest of this chapter covers only the one genuinely new class, \ccode{DualCoreDIP}, since everything else (\ccode{Camera}, \ccode{Vga}, \ccode{CacheAwareFrame}/\ccode{CacheAwareFramePool}, \ccode{Core1}, \ccode{Worker}, \ccode{DualCoreJob}, \ccode{System}) is unchanged from \nameref{chap:fpgasteel-dualcore-forwarding}.

\section{Prerequisites}
\label{sec:fpgasteel-dualcore-dip-prerequisites}

Same toolchain, \term{SD card} and \term{Eclipse}/\term{OpenOCD} debug setup as \nameref{chap:fpgasteel-dualcore-forwarding}, including the plain \repopath{boot.script.jtag_debug} script and the same two \term{CPU1}-related debug-configuration requirements. The \term{OV7670} camera adapter board caveat from \nameref{chap:fpgasteel-i2c-example} (Section~\ref{sec:fpgasteel-i2c-prerequisites}) applies here too.

\section{FPGA Hardware Involved}
\label{sec:fpgasteel-dualcore-dip-hardware}

Unchanged: see Section~\ref{sec:fpgasteel-cache-forwarding-hardware}. Unlike \nameref{chap:fpgasteel-dualcore-bringup} and \nameref{chap:fpgasteel-dualcore-forwarding}, this project actually drives the \term{LED} \term{PIO}: one bit per \ccode{DualCoreDIP::Result} value, read directly by \repopath{main0.cpp} after each processed frame.

\section{Operating principles}
\label{sec:fpgasteel-dualcore-dip-operating-principles}

\subsection{The pipeline, stage by stage}
\label{subsec:fpgasteel-dualcore-dip-pipeline}

\ccode{DualCoreDIP::process()} runs the nine parallel stages and two single-core stages from Section~\ref{sec:fpgasteel-dualcore-dip-about} in a fixed sequence, always in the same order, one synchronization point after another:

\begin{figure}[htbp]
 \centering
 \includesvg[width=\linewidth]{FPGAsteelDipPipeline}
 \caption{Every green stage is one \texttt{forkJoin()} (both cores), every blue stage runs on \term{CPU0} alone. The two early exits (\texttt{EMPTY}, \texttt{ON\_BORDER}) are what let an empty or unusable scene skip most of the pipeline instead of paying for all eleven stages every frame.}
 \label{fig:fpgasteel-dualcore-dip-pipeline}
\end{figure}

\subsection{The fork/join mechanism, reapplied}
\label{subsec:fpgasteel-dualcore-dip-forkjoin}

\ccode{DualCoreDIP::forkJoin()} is line-for-line the same sequence as \nameref{chap:fpgasteel-dualcore-forwarding}'s \ccode{DualCoreFormatConverter::start()} + \ccode{awaitDone()}: dispatch \ccode{m_worker1} to \term{CPU1}, run \ccode{m_worker0} synchronously on \term{CPU0}, then busy-poll \ccode{m_worker1.status} until it is \ccode{Completed} or \ccode{Error}. \ccode{process()} calls it nine times, reconfiguring \ccode{m_worker0}/\ccode{m_worker1} (source/destination pointers, dimensions, and a phase field selecting which of the nine operations to run) before each call. Nothing about the dispatch mechanism itself changes between one call and the next; only the data the two \ccode{Worker}s are pointed at does.

\subsection{The three build variants}
\label{subsec:fpgasteel-dualcore-dip-build-variants}

The three sibling projects differ in exactly two independent compile-time choices, applied to otherwise identical code (verified by diff: \repopath{main0.cpp} and every other class are byte-identical across all three; only \repopath{DualCoreDIP.{h,cpp}} and \repopath{Config.h}'s \ccode{ICACHE_ENABLE} change):

\begin{table}[htbp]
 \centering
 \begin{tabular}{@{}l l l@{}}
  \toprule
  \textbf{Project}\footnotemark & \textbf{Step implementations} & \textbf{L1 instruction cache} \\
  \midrule
  (none) & \term{ARM} \term{NEON} intrinsics & enabled \\
  \lightmidrule
  \repopath{_no_neon} & plain scalar C++, same algorithm & enabled \\
  \lightmidrule
  \repopath{_no_neon_no_icache} & plain scalar C++, same algorithm & disabled \\
  \bottomrule
 \end{tabular}
 \caption{The three build variants of this project.}
 \label{tab:fpgasteel-dualcore-dip-build-variants}
\end{table}
\footnotetext{All three share the \texttt{dualcore\_dip\_example} prefix; only the distinguishing suffix is shown here.}

\cmd{dualcore_dip_example}'s \ccode{step01}, \ccode{step03}, \ccode{step04}, \ccode{step05}, \ccode{step07}, \ccode{step08}, and \ccode{step09} (Figure~\ref{fig:fpgasteel-dualcore-dip-pipeline}'s seven \term{NEON}-eligible stages) use \texttt{arm\_neon.h}\footnote{Header shipped with the \term{ARM} toolchain (Paragraph~\ref{subsec:fpgasteel-install-paths}), not a separate download.} intrinsics to process 8 or 16 pixels per instruction; the \repopath{_no_neon} variant replaces every one of them with a plain byte-at-a-time loop computing the exact same result, with no other change anywhere in the codebase. \repopath{_no_neon_no_icache} is that same scalar build with \ccode{ICACHE_ENABLE} (\repopath{Config.h}) set to false instead of true, so \ccode{System::initCore()} skips enabling the \term{L1} instruction cache on either core.

Since \repopath{main0.cpp} already reports per-frame processing time on the hex display, building and running all three side by side turns the family into a small, controlled benchmark: the same algorithm, on the same hardware, with \term{NEON} and the instruction cache each toggled independently, one variable at a time.

\section{How it works}
\label{sec:fpgasteel-dualcore-dip-how-it-works}

The following snippet is taken from \repopath{main0.cpp}:

\begin{lstlisting}[language=C++]
sys->prepareCacheSubsystem();
Core1 & core1 = Core1::getInstance();
core1.start();
while (core1.getStatus() != Core1::Status::Ready
 && core1.getStatus() != Core1::Status::Error) {}
CacheAwareFramePool camPool(CAMERA_POOL_SIZE, FRAME_WIDTH, FRAME_HEIGHT,
 CAMERA_COLOR_SCHEME);
CacheAwareFramePool vgaPool(VGA_POOL_SIZE, FRAME_WIDTH, FRAME_HEIGHT,
 Frame::ColorScheme::RGB24);
vga->init(vgaPool);
pclkResetCtrl->deassertReset();
sys->delay(1000);
camera->init(camPool);
hex->showDashes();
DualCoreDIP dip(FRAME_WIDTH, FRAME_HEIGHT);
while (true) {
 if (camPool.available() > 0) {
  CacheAwareFrame * camFrame = camPool.pop(true);
  CacheAwareFrame * vgaFrame = static_cast<CacheAwareFrame *>(vgaPool.borrow());
  uint64_t t0 = sys->millis();
  DualCoreDIP::Result res = dip.process(*camFrame, *vgaFrame);
  uint64_t t1 = sys->millis();
  camPool.giveBack(camFrame);
  vgaPool.push(vgaFrame, true);
  LED_PIO = static_cast<uint8_t>(res);
  // LED0=OK, LED1=ON_BORDER, LED2=EMPTY
  hex->showNumber(static_cast<uint32_t>(t1 - t0));
 }
}
\end{lstlisting}

Structurally this is Section~\ref{sec:fpgasteel-dualcore-forwarding-how-it-works} with \ccode{dip.process(*camFrame, *vgaFrame)} standing in for \ccode{converter.convert(*camFrame, *vgaFrame)}, plus the one line writing the result code to \hw{LED_PIO}. The exception list, pool handling, and cache-flush placement around \ccode{pop(true)}/\ccode{push(vgaFrame, true)} are all identical.

\section{Debug tooling}
\label{sec:fpgasteel-dualcore-dip-debug-tooling}

\cmd{debugTools/view_frame.py}, same tool as Section~\ref{sec:fpgasteel-cache-forwarding-debug-tooling}. This project additionally ships \cmd{debugTools/linux_dualcore_emulator}, a small host-side (\term{Linux}, x86/\term{ARM}) build of \ccode{DualCoreDIP}'s pipeline logic against stub \term{hwlib} headers, used to test and tune the vision algorithm itself (thresholds, morphology, early-exit conditions) against saved frames on a development machine, without needing the board at all; it is not part of the bare-metal build and is not covered further here.